\section{The Problem Precedes the Mechanism}
The canonical question, ``How does the brain produce consciousness?'', appears mechanistic. It is actually downstream of a more basic question: \emph{what variable does the word consciousness denote?} Chalmers himself noted that ``consciousness'' is ambiguous before isolating conscious experience as the hard target \cite{Chalmers1995}. Block later called consciousness a ``mongrel concept'' and separated phenomenal from access consciousness \cite{Block1995}. Sanders et al. defined consciousness as subjective experience while distinguishing it from responsiveness and environmental connectedness \cite{Sanders2012}. Bayne, Hohwy, and Owen rejected a single levels-based description in favor of a multidimensional account \cite{Bayne2016}. LeDoux summarized the difficulty directly: what functions are assigned to consciousness depends on what one takes consciousness to be \cite{LeDoux2025}.

These are not merely verbal disagreements. A study of responsiveness, a study of dream experience, a study of memory, and a study of organismal viability can each be rigorous while targeting different variables. If all four are reported as measurements of one thing, the noun becomes part of the experimental design.

\defbox{\textbf{Target rule.} A mechanism can be evaluated only after its dependent variable is fixed. If the referent changes, apparent explanatory failure can be produced by target substitution rather than by the phenomenon.}

SCTC therefore makes an explicit revisionary assignment rather than allowing the noun to move between experiments.

\section{Definition Lock: What This Paper Means by Consciousness}
\subsection{The biological individual and its continuity}
Let $B$ denote a biological individual. Let
\begin{equation}
\Gamma_B(t_0,t_1)=1
\end{equation}
mean that the individual at $t_1$ lies on the continuing organism-level causal trajectory of the individual at $t_0$. Material turnover, repair, component replacement, architectural reorganization, and external support are compatible with continuity when they preserve the same biological individual.

Continuity across times and viability at a time are different objects. Define the living-system state
\begin{equation}
L_B(t)=
\begin{cases}
1, & B \text{ exists as a viable integrated biological individual at }t,\\
0, & B \text{ has reached terminal organism-level discontinuity.}
\end{cases}
\end{equation}
SCTC then fixes the minimal system-level target by definition:
\begin{equation}
\boxed{C_B(t)\equiv L_B(t).}
\label{eq:def}
\end{equation}
Equation \ref{eq:def} is primitive within the theory. Later cases do not prove it. The logical direction is
\begin{equation}
\text{definition}\;\rightarrow\;\text{deductions}\;\rightarrow\;\text{comparative tests}.
\end{equation}

\subsection{Faculties are not aliases for the system}
Let the declared faculty/state vector be
\begin{equation}
\mathbf F_B^*(t)=(W,E,K,R,M,P,Q,S,I),
\end{equation}
where $W$ is wakefulness/arousal; $E$ phenomenal experience; $K$ cognitive or experiential connectedness to the environment; $R$ behavioral responsiveness; $M$ memory; $P$ perceptual/sensory capacities; $Q$ cognitive access; $S$ self-awareness/metacognition; and $I$ intelligence or related cognitive capacities. Let $\mathcal A_B(t)$ denote the architecture supporting these capacities and $X_B(t)$ the state currently instantiated.

The hierarchy is descriptive, not set-theoretic:
\begin{equation}
B\text{ continuing}\;\longrightarrow\;\mathcal A_B\;\longrightarrow\;\mathbf F_B^*\;\longrightarrow\;X_B.
\end{equation}

\defbox{\textbf{Definition lock.} Throughout this paper, $C$ means system-level consciousness as defined in Eq. \ref{eq:def}. $E$ means phenomenal experience. No argument in this paper uses $C$ as a synonym for $E$. An objection that silently substitutes $E$ for $C$ is not an internal objection to SCTC; it is a rival target assignment and must be evaluated as such.}

A resolution rule is also required. $R$ and $K$ in $\mathbf F_B^*$ are behavioral/cognitive variables. Minimal biological responsiveness or metabolic coupling used to assess viability belongs to the implementation of $L_B$, not to the behavioral variables $R$ and $K$. Identical words at different biological scales do not denote identical variables.

\section{The Target-Substitution Error}
The most persistent misunderstanding can be stated formally.

\textbf{Principle 1 (Faculty exclusion).} If a declared faculty $F_j$ can be absent while $L_B(t)=1$, then $F_j$ is not constitutive of $C_B$ under SCTC:
\begin{equation}
L_B(t)=1\;\land\;F_j(t)=0 \;\Rightarrow\; F_j\not\equiv C_B.
\end{equation}
This is a consequence of the definition, not evidence that generated it.

\textbf{Principle 2 (Explanatory separation).} For any nonconstitutive faculty $F_j$,
\begin{equation}
F_j\not\equiv C_B
\quad\Rightarrow\quad
\mathrm{Unexplained}(F_j)\not\Rightarrow\mathrm{Unexplained}(C_B).
\label{eq:sep}
\end{equation}
Equation \ref{eq:sep} is the firewall between the theory and a recurring category error. If experience is not constitutive of system-level consciousness, then failing to explain why a particular physical process was experienced cannot by itself refute the definition of consciousness.

\subsection{The anesthesia dream, without the fog}
Consider an anesthetized person whose organismal continuity is preserved. Three epistemic cases are possible:
\begin{align}
L_B=1,\;R=0,\;E=1 &\quad \text{experience occurred},\\
L_B=1,\;R=0,\;E=0 &\quad \text{no experience occurred},\\
L_B=1,\;R=0,\;E=? &\quad \text{experience is unknown}.
\end{align}
Under SCTC, all three satisfy $C_B=1$. The scientifically valid question ``Why did a dream occur in the first case?'' targets $E$. It may be extremely difficult. It does not become a counterexample to Eq. \ref{eq:def} merely because the historical vocabulary called $E$ ``consciousness.''

This distinction can be summarized in one sentence:

\defbox{\textbf{A hard question about a faculty does not become a hard question about the bearer merely because both were given the same noun.}}

\section{Empirical Dissociations: What the Data Actually Separate}
The dissociation literature does not establish Eq. \ref{eq:def}; it demonstrates why the variables packed into the conventional noun should not be treated as interchangeable.

\subsection{Responsiveness, connectedness, and experience}
Sanders et al. explicitly argued that unresponsiveness is not equivalent to unconsciousness and distinguished subjective experience from responsiveness and connectedness \cite{Sanders2012}. Noreika et al. induced unresponsiveness with sedative/anesthetic agents and found later reports of subjective experience in almost 60\% of sessions \cite{Noreika2011}. Casey et al. evaluated EEG measures against explicit definitions of responsiveness, connectedness, and subjective experience, finding that measures derived from unresponsiveness did not map uniquely onto consciousness alone \cite{Casey2024}. Therefore
\begin{equation}
R=0\not\Rightarrow E=0.
\end{equation}
SCTC adds only the system-level statement that ordinary reversible anesthesia is compatible with $L_B=1$ while the faculty vector changes.

\subsection{Observable response and cognitive capacity}
Bodien et al. studied 353 adults with disorders of consciousness in a prospective, multicenter cohort. Among 241 participants without an observable response to commands, 60 (25\%) showed task-based command following on fMRI or EEG \cite{Bodien2024}. Thus
\begin{equation}
R_{\mathrm{observable}}=0\not\Rightarrow Q_{\mathrm{command}}=0.
\end{equation}
The scanner did not create the capacity. It exposed a capacity that the behavioral channel had failed to reveal.

\subsection{Memory and sensation}
Severe amnesia can coexist with preserved perception, language, reasoning, and continued biological existence \cite{Milner2005}. Individual sensory modalities can likewise be lost while the organism persists. These are elementary but decisive separations: memory, vision, hearing, or reportability may characterize the state of a conscious system without constituting the system-level category itself.

\begin{table*}[t]
\centering
\caption{The same continuing individual can occupy radically different faculty states. The table is not evidence for the primitive definition; it shows why the variables should not be used as synonyms.}
\footnotesize
\begin{tabularx}{0.96\textwidth}{@{}l c c c c X@{}}
\toprule
Condition & $L_B$ & $R$ & $E$ & $M$ & What is separated \\
\midrule
Ordinary wakefulness & 1 & often 1 & often 1 & variable & baseline faculty configuration \\
Reversible anesthesia & 1 & 0 & 0/1/? & often impaired & system continuity from response, experience, recall \\
Dreaming sleep & 1 & low & often 1 & variable & experience from environmental response \\
Severe amnesia & 1 & variable & possible & severely impaired & memory from system continuity \\
Sensory loss & 1 & variable & modality-specific loss & variable & modality from system continuity \\
Cognitive motor dissociation & 1 & 0 by behavior & unknown/possible & variable & observable response from cognitive command following \\
\bottomrule
\end{tabularx}
\end{table*}

\section{Why the ``Hard Problem of Consciousness'' Is Not a Neutral Phrase}
Chalmers' hard problem is the problem of why physical processing is accompanied by conscious experience \cite{Chalmers1995}. That is a coherent target. Under the notation here, it is the problem of $E$.

The phrase \emph{hard problem of consciousness} is therefore not definition-free. It becomes neutral only if $C\equiv E$ has already been established. Otherwise the phrase imports a contested identity into the name of the problem.

Under SCTC the correct question is
\begin{equation}
\boxed{\text{Which physical organizations instantiate }E,\text{ and why?}}
\end{equation}
That may remain the deepest problem in the science of mind. SCTC does not answer it by renaming it. SCTC denies that its unresolved status licenses the inference that system-level consciousness is likewise undefined.

The distinction is exactly the one exposed by the anesthesia example. Asking why a dream was experienced is a question about $E$. Asking whether the biological individual continued is a question about $L_B$. They are not two phrasings of the same dependent variable.

